% GENERATED FILE. Edit canonical lessons, then run npm run generate:books.
% canonical-source-hash: fnv1a64:2eec6ee301bb0978
% canonical-lessons: TA-C202-vara-iruti, TA-C202-varai, TA-C202-tiruvila, TA-C202-tipavali, TA-C202-puttantu, TA-R202-first-pass-words-from-chapters-194-197, TA-R202-second-pass-words-from-chapters-198-202

\chapter{Weekend, Up to, Temple festival, Deepavali, New Year}
\label{ch:ta-weekend-up-to-temple-festival-deepavali-new-year}
\hlchaptermodality{202}

\begin{chapteropening}
\textbf{By the end of this chapter:} \emph{I can say a weekend, up to, a temple festival, Deepavali and the New Year.}

Complete the last lesson of chapter 202: I can say a weekend, up to, a temple festival, Deepavali and the New Year.
\end{chapteropening}

\section[\texorpdfstring{\ta{வார} \ta{இறுதி} (vāra iṟuti)}{வார இறுதி (vāra iṟuti)}]{\ta{வார} \ta{இறுதி} (vāra iṟuti) — a weekend}

\label{lesson:TA-C202-vara-iruti}



\begin{quote}
Before the new one: say the Tamil for the palm of the hand, then the Tamil for a stomach ache.
\end{quote}

\subsection*{You'll want to know: \ta{வார} \ta{இறுதி}}
\textbf{\ta{வார} \ta{இறுதி}} — \emph{vāra iṟuti} — \textquotedblleft{}a weekend\textquotedblright{}.

\begin{grammarlens}[title={What you've built}]
One more word for this chapter.
\end{grammarlens}

\subsection*{Your turn}
\begin{itemize}
\raggedright
  \item \emph{Say it:} \emph{vāra iṟuti}
  \item \emph{Say it:} \emph{vāra iṟuti}, once more
  \item \emph{Say it:} say \emph{vāra iṟuti}
  \item \emph{Recall:} say the Tamil for the palm of the hand, then the Tamil for a stomach ache, then say \emph{vāra iṟuti} again
\end{itemize}

\begin{tcolorbox}[breakable,colback=teal!4,colframe=teal!35!black,title={Before you move on}]
What does \ta{வார} \ta{இறுதி} mean? (\textquotedblleft{}A weekend\textquotedblright{}.) Say it once more.
\end{tcolorbox}
\section[\texorpdfstring{\ta{வரை} (varai)}{வரை (varai)}]{\ta{வரை} (varai) — up to, until}

\label{lesson:TA-C202-varai}



\begin{quote}
Before the new one: say the Tamil for a stomach ache, then the Tamil for a weekend.
\end{quote}

\subsection*{You'll want to know: \ta{வரை}}
\textbf{\ta{வரை}} — \emph{varai} — \textquotedblleft{}up to, until\textquotedblright{}.

\begin{grammarlens}[title={What you've built}]
One more word for this chapter.
\end{grammarlens}

\subsection*{Your turn}
\begin{itemize}
\raggedright
  \item \emph{Say it:} \emph{varai}
  \item \emph{Say it:} \emph{varai}, once more
  \item \emph{Say it:} say \emph{varai}
  \item \emph{Recall:} say the Tamil for a stomach ache, then the Tamil for a weekend, then say \emph{varai} again
\end{itemize}

\begin{tcolorbox}[breakable,colback=teal!4,colframe=teal!35!black,title={Before you move on}]
What does \ta{வரை} mean? (\textquotedblleft{}Up to, until\textquotedblright{}.) Say it once more.
\end{tcolorbox}
\section[\texorpdfstring{\ta{திருவிழா} (tiruviḻā)}{திருவிழா (tiruviḻā)}]{\ta{திருவிழா} (tiruviḻā) — a temple festival}

\label{lesson:TA-C202-tiruvila}



\begin{quote}
Before the new one: say the Tamil for a weekend, then the Tamil for up to.
\end{quote}

\subsection*{You'll want to know: \ta{திருவிழா}}
\textbf{\ta{திருவிழா}} — \emph{tiruviḻā} — \textquotedblleft{}a temple festival\textquotedblright{}.

\begin{grammarlens}[title={What you've built}]
One more word for this chapter.
\end{grammarlens}

\subsection*{Your turn}
\begin{itemize}
\raggedright
  \item \emph{Say it:} \emph{tiruviḻā}
  \item \emph{Say it:} \emph{tiruviḻā}, once more
  \item \emph{Say it:} say \emph{tiruviḻā}
  \item \emph{Recall:} say the Tamil for a weekend, then the Tamil for up to, then say \emph{tiruviḻā} again
\end{itemize}

\begin{tcolorbox}[breakable,colback=teal!4,colframe=teal!35!black,title={Before you move on}]
What does \ta{திருவிழா} mean? (\textquotedblleft{}A temple festival\textquotedblright{}.) Say it once more.
\end{tcolorbox}
\section[\texorpdfstring{\ta{தீபாவளி} (tīpāvaḷi)}{தீபாவளி (tīpāvaḷi)}]{\ta{தீபாவளி} (tīpāvaḷi) — Deepavali, the festival of lights}

\label{lesson:TA-C202-tipavali}



\begin{quote}
Before the new one: say the Tamil for up to, then the Tamil for a temple festival.
\end{quote}

\subsection*{You'll want to know: \ta{தீபாவளி}}
\textbf{\ta{தீபாவளி}} — \emph{tīpāvaḷi} — \textquotedblleft{}Deepavali, the festival of lights\textquotedblright{}.

\begin{grammarlens}[title={What you've built}]
One more word for this chapter.
\end{grammarlens}

\subsection*{Your turn}
\begin{itemize}
\raggedright
  \item \emph{Say it:} \emph{tīpāvaḷi}
  \item \emph{Say it:} \emph{tīpāvaḷi}, once more
  \item \emph{Say it:} say \emph{tīpāvaḷi}
  \item \emph{Recall:} say the Tamil for up to, then the Tamil for a temple festival, then say \emph{tīpāvaḷi} again
\end{itemize}

\begin{tcolorbox}[breakable,colback=teal!4,colframe=teal!35!black,title={Before you move on}]
What does \ta{தீபாவளி} mean? (\textquotedblleft{}Deepavali, the festival of lights\textquotedblright{}.) Say it once more.
\end{tcolorbox}
\section[\texorpdfstring{\ta{புத்தாண்டு} (puttāṇṭu)}{புத்தாண்டு (puttāṇṭu)}]{\ta{புத்தாண்டு} (puttāṇṭu) — the New Year}

\label{lesson:TA-C202-puttantu}



\begin{quote}
Before the new one: say the Tamil for a temple festival, then the Tamil for Deepavali.
\end{quote}

\subsection*{You'll want to know: \ta{புத்தாண்டு}}
\textbf{\ta{புத்தாண்டு}} — \emph{puttāṇṭu} — \textquotedblleft{}the New Year\textquotedblright{}.

\begin{grammarlens}[title={What you've built}]
One more word for this chapter.
\end{grammarlens}

\subsection*{Your turn}
\begin{itemize}
\raggedright
  \item \emph{Say it:} \emph{puttāṇṭu}
  \item \emph{Say it:} \emph{puttāṇṭu}, once more
  \item \emph{Say it:} say \emph{puttāṇṭu}
  \item \emph{Recall:} say the Tamil for a temple festival, then the Tamil for Deepavali, then say \emph{puttāṇṭu} again
\end{itemize}

\begin{tcolorbox}[breakable,colback=teal!4,colframe=teal!35!black,title={Before you move on}]
What does \ta{புத்தாண்டு} mean? (\textquotedblleft{}The New Year\textquotedblright{}.) Say it once more.
\end{tcolorbox}
\section[(dialogue)]{Words from chapters 194-197 — a first pass}

\label{lesson:TA-R202-first-pass-words-from-chapters-194-197}



\begin{quote}
Say the words for Deepavali and the New Year.
\end{quote}

\subsection*{Your turn}
Say these aloud:

\begin{itemize}
\raggedright
  \item to pour, to chop, to peel, to deep-fry, to knead
  \item to bring, to give back, to drive, to use, to try
  \item to tear, to lose, to press, to feel, to soak
  \item to remind, to earn, to avoid, to move house, to contact
\end{itemize}

\begin{tcolorbox}[breakable,colback=teal!4,colframe=teal!35!black,title={Before you move on}]
To pour? (\textbf{\ta{ஊற்று}}, \emph{ūṟṟu}.) To chop? (\textbf{\ta{நறுக்கு}}, \emph{naṟukku}.) To peel? (\textbf{\ta{உரி}}, \emph{uri}.) To deep-fry? (\textbf{\ta{பொரி}}, \emph{pori}.) To knead? (\textbf{\ta{பிசை}}, \emph{picai}.) To bring? (\textbf{\ta{கொண்டு} \ta{வா}}, \emph{koṇṭu vā}.) To give back? (\textbf{\ta{திருப்பிக்} \ta{கொடு}}, \emph{tiruppik koṭu}.) To drive? (\textbf{\ta{ஓட்டு}}, \emph{ōṭṭu}.) To use? (\textbf{\ta{பயன்படுத்து}}, \emph{payaṉpaṭuttu}.) To try? (\textbf{\ta{முயற்சி} \ta{செய்}}, \emph{muyaṟci cey}.) To tear? (\textbf{\ta{கிழி}}, \emph{kiḻi}.) To lose? (\textbf{\ta{இழ}}, \emph{iḻa}.) To press? (\textbf{\ta{அழுத்து}}, \emph{aḻuttu}.) To feel? (\textbf{\ta{உணர்}}, \emph{uṇar}.) To soak? (\textbf{\ta{ஊற} \ta{வை}}, \emph{ūṟa vai}.) To remind? (\textbf{\ta{நினைவூட்டு}}, \emph{niṉaivūṭṭu}.) To earn? (\textbf{\ta{சம்பாதி}}, \emph{campāti}.) To avoid? (\textbf{\ta{தவிர்}}, \emph{tavir}.) To move house? (\textbf{\ta{குடிபெயர்}}, \emph{kuṭipeyar}.) To contact? (\textbf{\ta{தொடர்பு} \ta{கொள்}}, \emph{toṭarpu koḷ}.)
\end{tcolorbox}
\section[(dialogue)]{Words from chapters 198-202 — a second pass}

\label{lesson:TA-R202-second-pass-words-from-chapters-198-202}



\begin{quote}
Say the words for to pray and to cheat.
\end{quote}

\subsection*{Your turn}
Say these aloud:

\begin{itemize}
\raggedright
  \item to pray, to cheat, to tell a lie, to quarrel, to tease
  \item to believe, to agree, to refuse, to discuss, to compare
  \item to pinch, to bend, to paint, to collect, to suck
  \item the forehead, a beard, a moustache, the palm of the hand, a stomach ache
  \item a weekend, up to, a temple festival, Deepavali, the New Year
\end{itemize}

\begin{tcolorbox}[breakable,colback=teal!4,colframe=teal!35!black,title={Before you move on}]
To pray? (\textbf{\ta{பிரார்த்தனை} \ta{செய்}}, \emph{pirārttaṉai cey}.) To cheat? (\textbf{\ta{ஏமாற்று}}, \emph{ēmāṟṟu}.) To tell a lie? (\textbf{\ta{பொய்} \ta{சொல்}}, \emph{poy col}.) To quarrel? (\textbf{\ta{சண்டை} \ta{போடு}}, \emph{caṇṭai pōṭu}.) To tease? (\textbf{\ta{கேலி} \ta{செய்}}, \emph{kēli cey}.) To believe? (\textbf{\ta{நம்பு}}, \emph{nampu}.) To agree? (\textbf{\ta{ஒத்துக்கொள்}}, \emph{ottukkoḷ}.) To refuse? (\textbf{\ta{மறு}}, \emph{maṟu}.) To discuss? (\textbf{\ta{விவாதி}}, \emph{vivāti}.) To compare? (\textbf{\ta{ஒப்பிடு}}, \emph{oppiṭu}.) To pinch? (\textbf{\ta{கிள்ளு}}, \emph{kiḷḷu}.) To bend? (\textbf{\ta{வளை}}, \emph{vaḷai}.) To paint? (\textbf{\ta{பூசு}}, \emph{pūcu}.) To collect? (\textbf{\ta{சேகரி}}, \emph{cēkari}.) To suck? (\textbf{\ta{உறிஞ்சு}}, \emph{uṟiñcu}.) The forehead? (\textbf{\ta{நெற்றி}}, \emph{neṟṟi}.) A beard? (\textbf{\ta{தாடி}}, \emph{tāṭi}.) A moustache? (\textbf{\ta{மீசை}}, \emph{mīcai}.) The palm of the hand? (\textbf{\ta{உள்ளங்கை}}, \emph{uḷḷaṅkai}.) A stomach ache? (\textbf{\ta{வயிற்றுவலி}}, \emph{vayiṟṟuvali}.) A weekend? (\textbf{\ta{வார} \ta{இறுதி}}, \emph{vāra iṟuti}.) Up to? (\textbf{\ta{வரை}}, \emph{varai}.) A temple festival? (\textbf{\ta{திருவிழா}}, \emph{tiruviḻā}.) Deepavali? (\textbf{\ta{தீபாவளி}}, \emph{tīpāvaḷi}.) The New Year? (\textbf{\ta{புத்தாண்டு}}, \emph{puttāṇṭu}.)
\end{tcolorbox}
